\begin{lstlisting}[language=Python,caption={Ejercicio resuelto 2: FFT, IFFT y multiplicación}]
def resuelto_2():
    a, b = [1, 2, 3], [4, 5]
    padded = a + [0]
    roundtrip = ifft(fft(padded))
    result = multiply_fft(a, b)
    assert result == [4, 13, 22, 15]
    assert max(abs(x - y) for x, y in zip(padded, roundtrip)) < 1e-7
    return {"fft": fft(padded), "ifft_fft": roundtrip,
            "producto": result, "error_imaginario": max(abs(z.imag) for z in
             convolution_fft_raw(a, b))}
\end{lstlisting}
\noindent\textbf{Resultado de la ejecución:}
\begin{lstlisting}[style=consolestyle]
>>> resuelto_2()
{'fft': [(6+0j), (-1.9999999999999998+2j), (2+0j), (-2-2j)],
 'ifft_fft': [(1-0j),
              (2+5.721188726109833e-18j),
              (3-0j),
              -5.721188726109833e-18j],
 'producto': [4, 13, 22, 15],
 'error_imaginario': 6.661338147750939e-16}
\end{lstlisting}
